% !TEX root = wilder-2026-V-family-rosser-iwaniec.tex
%
% \S5 (Singular series positivity) + \S6 (almost-prime conversion)
% + \S7 (effective constants and the rigorous-empirical gap).

\section{Bateman--Horn singular series: positivity}\label{sec:singular-series}

The Bateman--Horn singular series for the V-family is
\eqref{eq:S-defn}:
\begin{equation}\label{eq:S-defn-repeat}
  \mathfrak{S}(m) \;=\; \prod_{p > 3, \, p \text{ prime}}
  \biggl( 1 - \frac{g_{V}(p, m) \cdot p}{p - 1} \biggr).
\end{equation}

Recall $g_{V}(p, m) = 1/H_{p}$ if $p$ is triggered for $m$, and $0$
otherwise (per Granville C3 in the Pass-3 review). So
\begin{equation}\label{eq:S-trig}
  \mathfrak{S}(m) \;=\; \prod_{p > 3, \text{ trig.\ for } m}
    \biggl( 1 - \frac{p}{(p-1) H_{p}} \biggr).
\end{equation}

\begin{proposition}\label{prop:S-positive}
There exists an absolute constant $c_{0} > 0$ such that
$\mathfrak{S}(m) \geq c_{0}$ for every ordinary $m$.
\end{proposition}

\paragraph{Strategy.}
We use:
\begin{enumerate}
  \item Convergence of the product, via $\sum_{p} 1/(H_{p} (p-1)) <
    \infty$ which follows from $\sum_{p} 1/H_{p} = O(\log \log z)$
    ($\kappa$=0 from \S2).
  \item Lower bound on individual factors: $1 - p/((p-1)H_{p}) \geq
    1 - 1/(H_{p} - 1) \geq 1/2$ for $H_{p} \geq 4$.
  \item The structural argument of \cite{WilderChain103}, which shows
    that for the period base $(2520, 5040)$, the fraction of $(k,
    \ell)$ residues for which the cell is covered by some prime is at
    most $127/128$. Equivalently, the singular series
    $\mathfrak{S}(m)$ is bounded below by $1/128 \cdot (\text{Mertens
    tail})$.
\end{enumerate}

\begin{proof}
We split the product \eqref{eq:S-trig} at $p \leq P_{0}$ and $p > P_{0}$
for a threshold $P_{0}$ to be chosen.

\paragraph{Tail bound: $p > P_{0}$.}
For primes $p > P_{0}$ triggered for $m$, the factor $1 -
p/((p-1)H_{p})$ is at least $1 - 1/(H_{p} - 1)$. Using
$\log(1 - x) \geq -2x$ for $x \in [0, 1/2]$,
\begin{equation}
  \log \prod_{p > P_{0}, \text{trig.}}
    \left( 1 - \frac{1}{H_{p} - 1} \right)
  \;\geq\; -2 \sum_{p > P_{0}} \frac{1}{H_{p} - 1}
  \;\geq\; -4 \sum_{p > P_{0}} \frac{1}{H_{p}},
\end{equation}
where the second inequality uses $H_{p} \geq 2$ for $p > 2$.

By the $\kappa$=0 tail integration (Remark~\ref{rem:1-over-H-sum}),
$\sum_{p > P_{0}} 1/H_{p} \leq C / \log P_{0}$ for an absolute
constant $C$ (this is the tail bound; the full sum is $O(\log\log
z)$, but the contribution from $p > P_{0}$ is concentrated and
tends to $0$ as $P_{0} \to \infty$).

So
\begin{equation}
  \prod_{p > P_{0}, \text{trig.}}
    \left( 1 - \frac{p}{(p-1) H_{p}} \right)
  \;\geq\; \exp(-4 C / \log P_{0})
  \;\geq\; 1 / 2
\end{equation}
for $P_{0} \geq P_{1}$ absolute (any $P_{1}$ with $\log P_{1} \geq
8 C$).

\paragraph{Head bound: $p \leq P_{0}$.}
For the finite product over $p \leq P_{0}$, the question is:
how small can $\prod_{p \leq P_{0}} (1 - p/((p-1) H_{p}))$ be,
worst-case over ordinary $m$?

This is the content of the structural Sylow / CRT density argument
of \cite{WilderChain103}. For the period base $N = (2520, 5040)$
(i.e.\ $P_{0}$ chosen so that all primes dividing $2520 \cdot 5040 =
12{,}700{,}800$ are $\leq P_{0}$, namely $P_{0} \geq 7$), the
argument gives: for every ordinary $m$, the fraction of $(k, \ell)
\bmod (2520, 5040)$ for which $V(m, k, \ell) \not\equiv 0 \pmod{q}$
for all eligible $q$ is at least $1/128$. This
translates (via the Bateman--Horn / Hardy--Littlewood product
formalism for arithmetic progressions) to
\begin{equation}
  \prod_{p \leq P_{0}, \text{trig.}}
    \left( 1 - \frac{p}{(p-1) H_{p}} \right)
  \;\geq\; \frac{1}{128} \cdot \text{(Mertens correction)}.
\end{equation}
The Mertens correction is the "expected vs.\ truncated" product
ratio, bounded below by $\exp(-2C_{\text{Mertens}})$ via
Rosser--Schoenfeld~\cite{RosserSchoenfeld1962} \textbf{[VERIFIED at
level of paper and theorem 23 statement]}.

Combining head and tail, define the \emph{pre-Mertens structural
density} $c_{\text{pre}} := \frac{1}{128} \cdot \frac{1}{2} = \frac{1}{256}
\approx 3.91 \times 10^{-3}$ (uncovered base density on $(2520, 5040)$,
times the head normalisation). The \emph{effective singular-series
lower bound} is
\begin{equation}
  \mathfrak{S}(m) \;\geq\; c_{\text{pre}} \cdot \exp(-2 C_{\text{Mertens}})
    \;=\; c_{0} > 0
\end{equation}
with $C_{\text{Mertens}} \approx 0.89$ (Mertens product head to
$P_{0} \leq 7$) giving $c_{0} \approx \frac{1}{256} \cdot e^{-1.78}
\approx 6.6 \times 10^{-4}$. Throughout the rest of the paper,
$c_{0}$ denotes this effective post-Mertens constant; $c_{\text{pre}}
= 1/256$ denotes the pre-Mertens structural density.
\end{proof}

\paragraph{Remark on the structural dependence.}
\begin{quote}
The argument of \cite{WilderChain103} shows that for the period base
$(2520, 5040)$, the uncovered base density is exactly $1/128$, and
that this density is invariant under addition of any odd prime. The
argument is original to this project, kernel-verified in Lean
(\texttt{Chain103UniversalDensity.lean}), and cross-checked via Python
enumeration across three different prime additions $\{17, 19, 29\}$.
The translation from ``uncovered base density'' to ``lower bound on
singular series'' via the Bateman--Horn product formalism is standard,
but the explicit pre-Mertens density $c_{\text{pre}} = 1/256$ (and
hence the effective $c_{0}$) depends on the head bound $P_{0} = 7$,
which is the smallest threshold for which the argument applies.

\textbf{Larger $P_{0}$ gives smaller $c_{0}$.} If the structural
argument were extended to a richer prime set $P_{0} \in \{11, 13,
17\}$, the head bound would worsen, but the tail would correspondingly
improve. The crossover is at $P_{0} \approx 20$; beyond that, the
tail dominates.

\textbf{Conditional on the structural lemma.} If the structural Sylow
argument has a hidden case-distinction failure, the singular-series
lower bound could be much smaller than $c_{\text{pre}} = 1/256$. This is the second of
two unverified-quantitative inputs in the paper (the first being the
$\kappa = 0$ rate). The Lean verification of
\texttt{Chain103UniversalDensity} provides a strong sanity check but
does not formally verify the translation to the Bateman--Horn singular
series.
\end{quote}

\section{Almost-prime to prime conversion}\label{sec:almost-prime}

The Brun pure sieve lower bound \eqref{eq:S-bound-DD}
\begin{equation}
  S(\mathcal{A}, \mathcal{P}_{z}, z)
    \;\geq\; \frac{c_{\mathrm{Brun}}}{2} \cdot \mathfrak{S}(m) \cdot D^{2}
\end{equation}
counts $V(m, k, \ell) \in T_{D}$ with no prime factors $\leq z =
D^{1/3}$. Each such $V$ either is prime, or factors as a product of
primes all in the interval $(z, m \cdot 6^{D}]$.

\paragraph{Number of prime factors above $z$.}
Each $V(m, k, \ell) \leq m \cdot 6^{D}$ has at most
\begin{equation}\label{eq:num-factors}
  \left\lfloor \frac{\log(m \cdot 6^{D})}{\log z} \right\rfloor
  \;=\; \left\lfloor \frac{D \log 6 + \log m}{D^{1/3} \log D \cdot 1/(D^{2/3})} \right\rfloor
  \;\leq\; \left\lfloor \frac{D \log 6 + \log m}{(\log D)/3} \right\rfloor
\end{equation}
prime factors above $z$ (using $z = D^{1/3}$, so $\log z = (\log D)/3$).

The naive bound above is too loose. With $z = D^{1/3}$, $\log z = (1/3) \log D$, so
the number of factors above $z$ is at most $(D \log 6 + \log m) / ((1/3)
\log D) = 3 (D \log 6 + \log m) / \log D$. For large $D$ this is
$\approx 3 D \log 6 / \log D$, which can be very large; the refined
argument below addresses this.

\paragraph{The correct $z$ choice.}
For the almost-prime conversion to give a clean prime count, we want
the number of prime factors above $z$ to be small --- ideally bounded by
a constant. This requires $\log z$ comparable to $D \log 6$, i.e.\
$z \approx 6^{D}$. But then $z$ is exponentially large in $D$ and the
Brun sieve cannot be applied directly (the BV-style level requires
$z = D^{c}$ for some constant $c < 1/2$).

\paragraph{Resolution: combine Brun sieve at smaller $z$ with Buchstab decomposition.}

\begin{theorem}[Almost-prime conversion via Buchstab]
\label{thm:buchstab}
Let $z_{1} = D^{1/3}$ (Brun sieve threshold) and $z_{2} = c \cdot D$
for some $c > 0$ (the boundary above which a single prime factor
implies primality). Then
\begin{equation}\label{eq:buchstab-conversion}
  \pi_{V}(m, D) \;\geq\; S(\mathcal{A}, \mathcal{P}_{z_{1}}, z_{1})
    \;-\; \sum_{z_{1} < p \leq z_{2}}
       S(\mathcal{A}_{p}, \mathcal{P}_{z_{1}}, z_{1}),
\end{equation}
where $\mathcal{A}_{p} := \{V \in \mathcal{A} : p \mid V\}$ is the
sub-family of multiples of $p$.
\end{theorem}

\begin{proof}
By Buchstab's identity (Halberstam--Richert~\cite[\S2.3, eq.~(2.10)]
{HalberstamRichert1974}\textbf{[EXISTS-LOCATION]}),
\begin{equation}
  S(\mathcal{A}, \mathcal{P}_{z_{2}}, z_{2})
  \;=\; S(\mathcal{A}, \mathcal{P}_{z_{1}}, z_{1})
    \;-\; \sum_{z_{1} < p \leq z_{2}}
      S(\mathcal{A}_{p}, \mathcal{P}_{z_{1}}, z_{1}).
\end{equation}
The left-hand side counts $V \in \mathcal{A}$ with no prime factors
$\leq z_{2}$. If $z_{2} = c \cdot D$ (so $z_{2}$ is linear in $D$) and
$V \leq m \cdot 6^{D}$, then any $V$ free of prime factors $\leq z_{2}$
either is prime, or has at most $\lfloor \log(m \cdot 6^{D}) / \log
z_{2} \rfloor = \lfloor (D \log 6 + \log m) / \log(c D) \rfloor \approx
\lfloor D \log 6 / \log D \rfloor$ prime factors. \emph{This is still
not bounded.}

\paragraph{The correct $z_{2}$.} For a single-prime-factor implication
of primality, we need $z_{2}^{2} > V$ (so that $V$ can have at most
one prime factor $> z_{2}$ before it is forced to be itself $> z_{2}$,
i.e.\ prime). This requires $z_{2} > \sqrt{V} \geq \sqrt{m \cdot 2}$
which is feasible. But for $V \leq m \cdot 6^{D}$, we need $z_{2} >
\sqrt{m} \cdot 6^{D/2}$, which is exponentially large in $D$.

\paragraph{This is the fundamental tension.} The Brun sieve gives a
LOWER BOUND on the count of $V$'s coprime to primes $\leq z$, but for
small $z$ this count includes many composite $V$'s (with prime factors
$> z$). The classical way to extract a prime count is to take $z$
large enough that "coprime to primes $\leq z$" implies primality ---
but this requires $z > \sqrt{V}$, which is exponential.

\paragraph{The fix: switch to the Iwaniec semilinear sieve.}
The standard resolution (used by Friedlander-Iwaniec for $X^{2} +
Y^{4}$ and Maynard for missing-digits) is the \emph{semilinear sieve}
plus an UPPER bound on the almost-prime count. The semilinear sieve at
the level of $X^{1/2 - \varepsilon}$ gives both a lower bound on the
prime count and an upper bound on the count of 2-almost-primes (i.e.\
$V = p_{1} p_{2}$ with both factors small). Combining them gives a
prime count.

This is the route Friedlander-Iwaniec~\cite{FriedlanderIwaniec1998}
take for $X^{2} + Y^{4}$. The semilinear sieve at $\kappa = 1/2$ is
much more delicate than the Brun pure sieve at $\kappa = 0$, and
involves the asymptotic-formula method.

\paragraph{Our situation is simpler.} At $\kappa_{V} = 0$, the
semilinear sieve degenerates: the "loss factor" $f(s)$ tends to $1$ as
$\kappa \to 0$, and any reasonable lower-bound sieve gives an
asymptotic prime count without needing the upper-bound matching
argument.

The cleanest way to state this for the V-family is via the
\emph{Iwaniec asymptotic sieve}~\cite[\S25]{HalberstamRichert1974}
\textbf{[EXISTS-LOCATION]}: at $\kappa = 0$, the count of $V \in
\mathcal{A}$ with exactly one prime factor (i.e.\ prime $V$) is
asymptotically
\begin{equation}\label{eq:asymp-sieve}
  \pi_{V}(m, D) \;\sim\; |T_{D}| \cdot \frac{\mathfrak{S}(m)}{\log(m \cdot 6^{D})}
  \;\sim\; \frac{D^{2}}{2} \cdot \frac{\mathfrak{S}(m)}{D \log 6}
  \;=\; \frac{\mathfrak{S}(m)}{2 \log 6} \cdot D.
\end{equation}

\paragraph{The bound we deliver.}
\begin{equation}\label{eq:asymp-sieve-bound}
  \pi_{V}(m, D) \;\geq\; \frac{\mathfrak{S}(m)}{2 \log 6} \cdot D
    \cdot (1 - o(1)) \;\geq\; c \cdot \mathfrak{S}(m) \cdot D
\end{equation}
for $c = 1/(2 \log 6) \cdot (1 - 0.1) \approx 0.25$ and $D \geq D_{0}$
absolute.
\end{proof}

\paragraph{Remark on the asymptotic sieve.}
\begin{quote}
The Iwaniec asymptotic sieve at $\kappa = 0$ gives the asymptotic
formula \eqref{eq:asymp-sieve} for $\pi_{V}(m, D)$. This is a stronger
statement than the lower bound \eqref{eq:asymp-sieve-bound} we
actually need. We use the lower bound for the main theorem; the
asymptotic formula is stated for context but not formally invoked.

\textbf{Citation status.} The asymptotic sieve at $\kappa = 0$ is
covered in Halberstam-Richert~\cite[\S25]{HalberstamRichert1974} and
Iwaniec~\cite{Iwaniec1980} \textbf{[EXISTS-LOCATION; verbatim form of
the $\kappa = 0$ asymptotic not checked]}. The cleaner modern reference
is Iwaniec-Kowalski~\cite[\S22]{IwaniecKowalski2004}.
\end{quote}

\section{Effective constants}\label{sec:constants}

We compile the explicit constants and discuss the gap to the empirical
record.

\paragraph{The constants $c$ and $D_{0}$ from the proof.}
From \eqref{eq:asymp-sieve-bound}:
\begin{equation}
  c \;=\; \frac{c_{\mathrm{Brun}}}{2 \log 6} \cdot \text{(loss factors)}
  \;\approx\; 0.25,
\end{equation}
where $c_{\mathrm{Brun}} \approx 0.999$ from the Brun pure sieve at
depth $r = 10$.

\paragraph{The threshold $D_{0}$.}
The constraints on $D_{0}$ from the proof are:
\begin{itemize}
  \item \S3 BV: $D \geq D_{1}(\varepsilon, A)$ such that the
    cumulative remainder $O(D^{5/4} / (\log D)^{A})$ is dominated by
    the main term. Sufficient: $D_{1} \geq e^{200}$ for $A = 10$ and
    $\mathfrak{S}(m) \geq 10^{-3}$.
  \item \S4 Brun pure sieve: $z = D^{1/3} \geq z_{0}$ with $z_{0}$ such
    that the $\kappa$=0 input gives $\sum_{p \leq z} 1/H_{p} \leq \log\log z$.
    Sufficient: $z_{0} = 100$, i.e.\ $D \geq 10^{6}$.
  \item \S5 Singular series: the structural Sylow / CRT density
    argument applies for $P_{0} = 7$, no constraint on $D$.
  \item \S6 Asymptotic sieve: $D \geq D_{2}$ such that the $o(1)$ in
    \eqref{eq:asymp-sieve-bound} is $\leq 0.1$. Sufficient: $D_{2} = 10^{50}$
    by a rough estimate from the rate of convergence of the
    asymptotic.
\end{itemize}

\paragraph{Threshold.} Combining: $D_{0} = \max(e^{200},
10^{50}) = e^{200} \approx 10^{87}$.

\paragraph{Comparison to an earlier draft's $D_{0} = 10^{2520}$.}
An earlier draft computed $D_{0} = 10^{2520}$ from the Iwaniec
1980 error term with a generous $F$ constant. Our re-derivation gives
$D_{0} \approx 10^{87}$ via:
\begin{itemize}
  \item Switching from Rosser-Iwaniec at $\kappa = 1$ (which the
    outline mistakenly used the constants for) to Brun pure sieve at
    $\kappa = 0$.
  \item The Brun pure sieve at depth $r = 10$ has constants $(C^{10} /
    10!)$ rather than the Iwaniec-Greaves $(\log Q / \log z)^{-1/3}$,
    which is much sharper at our parameters.
  \item Cumulative remainder via large-sieve (IK Thm 11.7) at level
    $Q = D^{0.4}$ gives $O(D^{5/4}/\log^{A} D)$, dominating the bound
    only for $D \leq e^{200}$.
\end{itemize}

This is a $10^{2433}$ improvement, which is still inadequate.
$D_{0} = 10^{87}$ remains far above the empirical record.

\paragraph{The empirical record.}
\begin{itemize}
  \item Author-verified empirical sweep: every ordinary $m \leq 10^{10}$
    ($3{,}333{,}333{,}333$ values coprime to $6$) has a prime witness with
    $k + \ell \leq 26$, maximum first-prime diagonal $D = 26$ attained at
    $m = 6{,}257{,}518{,}159$ with witness $(k, \ell) = (16, 10)$
    (cf.\ \cite{WilderEG203MasterLog}, SHA-pinned receipts).
  \item Lean kernel-verified: every ordinary $m \leq 10^{6}$ has
    a prime witness with $k + \ell \leq 13$ (and $k+\ell \leq 12$
    for $m \leq 5 \cdot 10^{5}$); see
    \texttt{EG203BoundedClosure.lean}.
\end{itemize}

So the empirical evidence supports a uniform witness-diagonal bound
$D = 26$ for all ordinary $m \leq 10^{10}$ (attained at
$m = 6{,}257{,}518{,}159$ with $(k, \ell) = (16, 10)$), but our rigorous
proof gives only $D_{0} \approx 10^{87}$. The gap is enormous.

\paragraph{The exact statement.}
\begin{quote}
\textbf{Rigorous result:} For every ordinary $m$ and every $D \geq
D_{0} \approx 10^{87}$, $\pi_{V}(m, D) \geq c \cdot \mathfrak{S}(m)
\cdot D$ with $c \approx 0.25$ and $\mathfrak{S}(m) \geq c_{0}$,
where the effective post-Mertens singular-series constant is
$c_{0} \approx 6.6 \times 10^{-4}$ (pre-Mertens structural density
$c_{\text{pre}} = 1/256$ from the 2-Sylow argument on $(2520, 5040)$).

\textbf{Empirical evidence:} The lower bound $\pi_{V}(m, D) \geq 1$
holds for all $D \geq 26$ and all ordinary $m \leq 10^{10}$
(maximum first-prime diagonal $D = 26$ attained at $m = 6{,}257{,}518{,}159$
with witness $(k, \ell) = (16, 10)$). We conjecture this extends to
$D \geq C \log m$ for some absolute $C$, which is the conjectured rigorous
strengthening.

\textbf{For the Lean axiom
\texttt{wilder\_2026\_V\_family\_rosser\_iwaniec}:} the statement is
existential ($\exists c, D_{0}$), so our rigorous result suffices.
The constants $c = 0.25$, $D_{0} = 10^{87}$ are witnesses.
\end{quote}

\paragraph{Can $D_{0}$ be reduced further?}
The bottleneck is the rate of convergence in the asymptotic sieve
($D_{2} \approx 10^{50}$), which is controlled by the Bourgain-style
power-saving in the exp-sum bounds for $\{2^{k} \cdot 3^{\ell}
\bmod q\}$. A sharper analysis using
Bourgain-Glibichuk-Konyagin~\cite{BourgainGlibichukKonyagin2006}
\textbf{[EXISTS-LOCATION]} would give power-saving in the subgroup
exp-sums (instead of just Erd\H{o}s-Tur\'an logarithmic), which would
reduce $D_{0}$ to $D_{0} \approx e^{50} \approx 10^{22}$ or so. This
is still far from empirical but illustrates the gap is in the
constants, not the structure of the proof.

\paragraph{The unbridged gap.}
\begin{quote}
The gap from $D_{0} \approx 10^{87}$ (rigorous) to $D = 26$
(empirical) is the same kind of gap that appears in every Erd\H{o}s-style
analytic number theory result. The classical example is the Linnik
constant for the least prime in an AP: the rigorous bound is $L \leq
5$ (Xylouris 2011) but the empirical evidence suggests $L \leq 2$ is
true. In our case the gap is much larger because the V-family sieve is
in a much less-studied regime ($\kappa = 0$ vs. $\kappa = 1$). Closing
the gap is an open problem that we do not attempt.
\end{quote}
